% =====================================================================
%  Seance 5 -- C. supervise : quand on a des etiquettes historiques
% =====================================================================
\section{Quand on a des étiquettes : le détecteur supervisé}

\begin{frame}[fragile]{Le scénario supervisé : les incidents passés sont des étiquettes}
\begin{columns}[t]
\begin{column}{0.5\textwidth}
\begin{block}{Etiqueter les sessions}
Ici, « vérité terrain » : les sessions de la fenêtre d'attaque réelle
(bdd-principal, 02 h -- 04 h le 05/10) $=$ \textbf{1}, les 9\,988 autres
$=$ \textbf{0}. En production, ces étiquettes viennent des tickets
d'incident, des rapports du SOC, des post-mortems.
\end{block}
\begin{itemize}
\item Problème : \textbf{12 positives sur 10\,000} --- déséquilibre extrême
(0,12 \%), comme les 120 fenêtres incident de la séance 2.
\item Réponses : \py{class\_weight="balanced"}, forêt robuste, évaluation
stratifiée.
\item Split : 50 \% train / 50 \% test \textbf{stratifié}, graine 5534 ---
6 positives en test.
\end{itemize}
\end{column}
\begin{column}{0.46\textwidth}
\begin{codebox}[Le cœur de B4]
\begin{lstlisting}
rf = RandomForestClassifier(
    n_estimators=300,
    class_weight="balanced",
    random_state=5534, n_jobs=2)
rf.fit(X_train, y_train)
pred = rf.predict(X_test)
precision_score(y_test, pred)
recall_score(y_test, pred)
\end{lstlisting}
\end{codebox}
\end{column}
\end{columns}
\end{frame}

\begin{frame}{Resultat réel : non-supervisé vs supervisé}
\begin{columns}[t]
\begin{column}{0.5\textwidth}
\begin{center}
\includegraphics[width=\textwidth]{f32_supervise_vs_non.png}
\end{center}
\end{column}
\begin{column}{0.46\textwidth}
\begin{ebox}[Ce que disent les chiffres]
\begin{itemize}
\item Supervisé : précision \textbf{1,00} et rappel \textbf{1,00} sur le
test (6 positives) --- sur ce jeu, le signal est si net que la forêt sépare
parfaitement.
\item Non supervisé : rappel 1,00 aussi, mais précision@30 $= 0{,}40$ ---
\emph{sans aucune étiquette}.
\end{itemize}
\end{ebox}
\begin{wbox}[Méfiez-vous du 100 \%]
6 positives en test : l'intervalle de confiance est énorme. Un modèle
parfait sur un jeu petit et propre échoue en production (attaque
\emph{différente}, données bruitées). Le 100 \% mesure le jeu, pas la
sécurité du parc.
\end{wbox}
\end{column}
\end{columns}
\end{frame}

\begin{frame}{Ce que la forêt a appris : importance des features}
\begin{center}
\footnotesize
\begin{tabular}{@{}lcc@{}}
\toprule
\textbf{Feature} & \textbf{Importance mesurée} & \textbf{Lecture}\\
\midrule
\py{comptes\_sensibles} & 0,259 & viser \py{root}/\py{admin}/\py{inconnu-*} : signal fort\\
\py{echecs} & 0,226 & le volume d'échecs, attendu\\
\py{echecs\_consecutifs} & 0,219 & la rafale ininterrompue\\
\py{taux\_echec} & 0,213 & zéro succès pendant l'attaque\\
\py{utilisateurs} & 0,051 & contribution secondaire\\
\py{entropie} & 0,030 & redondante avec les échecs ici\\
\py{nuit} & 0,002 & presque inutile sur ce jeu\\
\bottomrule
\end{tabular}
\end{center}
\vspace{1mm}
\begin{kbox}[Les quatre features de tête racontent la même histoire]
Comptes sensibles + échecs + rafale + zéro succès $=$ la \emph{définition}
d'un brute force. La forêt n'a rien inventé : elle a redécouvert la
signature --- et \textbf{quantifié} la contribution de chaque dimension,
ce qu'une règle écrite à la main ne fait pas.
\end{kbox}
\end{frame}

\begin{frame}{Choisir son détecteur : arbre de décision opérationnel}
\begin{columns}[t]
\begin{column}{0.52\textwidth}
\begin{block}{Vous n'avez AUCUNE étiquette historique}
$\Rightarrow$ non supervisé (IsolationForest, LOF) + validation par
échantillonnage manuel du top-N. C'est le point d'entrée de \emph{tout}
SOC qui démarre.
\end{block}
\begin{block}{Vous avez des incidents étiquetés, rares}
$\Rightarrow$ supervisé à \py{class\_weight} équilibré ; métriques :
\textbf{rappel d'abord} (rater une intrusion coûte cher), précision ensuite.
\end{block}
\end{column}
\begin{column}{0.44\textwidth}
\begin{block}{Vous avez les deux}
$\Rightarrow$ les deux en production (séance 5 d : redondance des
détecteurs) + le supervisé réentraîné à chaque incident confirmé : la
boucle d'amélioration.
\end{block}
\begin{qbox}[Question]
Vos 10\,000 sessions n'ont \emph{aucune} étiquette fiable. Pourquoi ne pas
étiquetter « session à $>$ 5 échecs $=$ attaque » et entraîner là-dessus ?\\
(a) c'est trop lent \quad (b) c'est circulaire : le modèle apprendrait votre
règle, jamais mieux qu'elle \quad (c) les forêts n'aiment pas les étiquettes
\end{qbox}
\end{column}
\end{columns}
\end{frame}

\begin{frame}{Quiz 3 --- supervisé et évaluation (5 min)}
\begin{qbox}[3 questions]
\begin{enumerate}
\item Avec 12 positives sur 10\,000 sessions, la métrique à surveiller en
priorité est :\\
(a) l'exactitude globale \quad (b) le rappel et la précision sur la classe
rare \quad (c) le score R²
\item Précision@30 $= 0{,}40$ signifie :\\
(a) 40 \% de l'attaque est bloquée \quad (b) parmi les 30 sessions les plus
anormales, 12 sont de vraies attaques \quad (c) le modèle se trompe 40 \% du temps
\item Le meilleur usage des incidents confirmés par les analystes est :\\
(a) les archiver \quad (b) en faire des étiquettes pour réentraîner le
détecteur supervisé \quad (c) baisser le seuil d'alerte
\end{enumerate}
\end{qbox}
\vspace{1mm}
\begin{solbox}[Réponses]
\textbf{1 : b} --- l'exactitude est piégeuse (99,88 \% en prédisant « rien »).
\quad
\textbf{2 : b} --- précision@N juge le \emph{classement}, l'outil de l'analyste.
\quad
\textbf{3 : b} --- c'est la boucle non supervisé $\rightarrow$ supervisé.
\end{solbox}
\end{frame}
